\section{Follow-up: can stronger dropout reduce overfitting?}
\label{sec:dropoutdesign}
\textbf{Why regularization, and which kind?} The rising validation losses motivated this follow-up; stronger dropout was not part of the initial model-size hypothesis. The baseline already uses dropout 0.1 and AdamW weight decay 0.1, so we investigate whether stronger regularization helps.

Two relevant options are weight decay and dropout. Weight decay shrinks weights during optimization, discouraging large weight values. AdamW separates this shrinkage from the adaptive gradient update \citep{loshchilov2019adamw}. Dropout instead randomly removes activations during training, reducing reliance on particular feature combinations \citep{srivastava2014dropout}. For drop probability $p$,
\begin{equation}
\widetilde h=\frac{m\odot h}{1-p},\qquad m_j\sim\mathrm{Bernoulli}(1-p).
\end{equation}
Here, $h$ is an activation and $m$ is a random zero/one mask. Remaining values are rescaled to keep their expected size unchanged. Dropout is disabled for evaluation. The shared probability affects embeddings, attention probabilities, attention outputs, and MLP outputs.

We chose to increase dropout from 0.1 to 0.2 while keeping weight decay unchanged. This lets us test one setting within the remaining budget. The readings motivate both options, but do not establish which is best for this dataset. We did not measure weight norms or test stronger weight decay, so we cannot claim that insufficient weight decay caused the problem or that dropout is the better remedy. The value 0.2 is a moderate trial, not a known optimum.

\textbf{Controlled comparison.} All three sizes are retrained from random initialization. Within each size, we keep the architecture, seed 1337, AdamW settings, learning-rate schedule, context 1,024, batch 32, validation windows, evaluation frequency, and 7,200-second budget unchanged. Only dropout and output location change. The launcher checks configurations against the saved baselines. Including small checks whether an intervention motivated by medium and large also helps a model with less obvious overfitting. This adds three two-hour runs, reusing the earlier batch-size benchmark.

\textbf{Success criterion.} Compare each size with its own baseline using the lowest recorded validation loss. A smaller training--validation gap alone is insufficient if validation loss worsens. Final loss, time of the best checkpoint, and tokens processed are secondary checks. One shared seed controls initialization, but does not make dropout masks or training paths identical. This is a limited follow-up, not an exhaustive hyperparameter search.
